% chapter: 付録（問いの連鎖）

\chapter{付録}\label{章：付録}
\begin{bookcolumns}

  付録では,本編で選んだ量や変形を別の表現から見直す.
  \begin{itemize}
    \item 連立型や隣接3項間をまとめて更新したいなら,行列の基本から読む.
    \item 階差・総和と微積分の関係を見たいなら,\sectionref*{節：差分から始まる離散的な微積分}へ進む.微分・積分を使う部分は任意に読み進める.
    \item 数列全体の係数を扱いたいなら,\sectionref*{節：数列全体を一つの対象にまとめると,何が見えるか}へ進む.母関数の係数計算は,最初から無限級数の収束を仮定する計算とは区別する.
  \end{itemize}
  最初の一巡ではすべてを証明し直す必要はない.
  「本編のどの操作が現れたか」「新しい表現で何が見えたか」を一つずつ確かめてほしい.

% 付録：行列の基本

\section{行列の基本}\label{節：行列の基本}

本編では,複雑な\bookterm{recurrence}{漸化式}を等比型などの既知の形へ変えることを繰り返してきた.
付録では,その置き換えがなぜうまくいくのかを,\bookterm{matrix}{行列}や\bookterm{difference-operator}{差分},\bookterm{generating}{母関数}という言葉で捉え直す.
最初に,そのための道具となる\bookterm{matrix}{行列}の読み方と計算の仕方を説明しよう.
\bookterm{matrix}{行列}を初めて学ぶ読者も,数値例から順に確かめてほしい.
ここでは主に二行二列の\bookterm{matrix}{行列}を扱う.
後の節では,本編で求めた\bookterm{fixed}{不動点}や\bookterm{geometric}{等比数列}が,ここで学ぶ計算の中にそのまま現れる.

\subsection{行列とは何か}
  \BookFigMatrixMap


\bookterm{matrix}{行列}とは,数を長方形に並べ,全体を括弧で囲んだものである.
たとえば
\[
  A=\begin{pmatrix}2&1\\1&2\end{pmatrix}
\]
は,縦に二段,横に二列の数を並べた\bookterm{matrix}{行列}である.
横の並びを\textbf{行},縦の並びを\textbf{列}という.
上が第一行,下が第二行であり,左が第一列,右が第二列である.
並んでいる一つ一つの数を\textbf{成分}と呼ぶ.
この例の第一行・第二列の成分は$1$である.

また,二つの数を縦に並べた
\[
  \boldsymbol{v}=\begin{pmatrix}x\\y\end{pmatrix}
\]
を\textbf{\bookterm{column-vector}{列ベクトル}}と呼ぶ.
太字の$\boldsymbol{v}$は,二つの数をまとめて表すための記号である.
平面上の点$(x,y)$への矢印と考えることもできるが,以下ではまず「二つの値を入れた一組の記録」として読めばよい.

\bookterm{column-vector}{列ベクトル}の足し算と定数倍は,成分ごとに行う.
\[
  \begin{pmatrix}x\\y\end{pmatrix}
  +\begin{pmatrix}u\\v\end{pmatrix}
  =\begin{pmatrix}x+u\\y+v\end{pmatrix},
  \qquad
  c\begin{pmatrix}x\\y\end{pmatrix}
  =\begin{pmatrix}cx\\cy\end{pmatrix}.
\]
たとえば$\begin{pmatrix}3\\1\end{pmatrix}+\begin{pmatrix}1\\-1\end{pmatrix}
=\begin{pmatrix}4\\0\end{pmatrix}$である.
二つの\bookterm{column-vector}{列ベクトル}が等しいとは,上下それぞれの成分が等しいことをいう.
一般の二成分のベクトルも
\[
  \begin{pmatrix}x\\y\end{pmatrix}
  =\frac{x+y}{2}\begin{pmatrix}1\\1\end{pmatrix}
   +\frac{x-y}{2}\begin{pmatrix}1\\-1\end{pmatrix}
\]
と分けられる.上下の成分をそれぞれ足して確かめてほしい.
本編の\sectionref*{節：連立型}で和と差を求めた操作は,この二つの係数を求めることに対応する.
この対応は\bookproblemref{practice-basic}{7}{標準問7}で確認できる.
後で,この二つの方向が更新の定数倍とどう結び付くかを見る.

\subsubsection{行列を列ベクトルに掛ける}

先ほどの\bookterm{matrix}{行列}$A$を\bookterm{column-vector}{列ベクトル}に掛けるときは,次のように計算する.
\[
  \begin{pmatrix}2&1\\1&2\end{pmatrix}
  \begin{pmatrix}x\\y\end{pmatrix}
  =
  \begin{pmatrix}2x+y\\x+2y\end{pmatrix}.
\]
上の成分は,第一行の$2,\,1$を$x,y$にそれぞれ掛けて足したもの.
下の成分は,第二行の$1,\,2$を$x,y$にそれぞれ掛けて足したものである.
たとえば
\[
  A\begin{pmatrix}1\\0\end{pmatrix}
  =\begin{pmatrix}2\\1\end{pmatrix},
  \qquad
  A\begin{pmatrix}2\\1\end{pmatrix}
  =\begin{pmatrix}5\\4\end{pmatrix}.
\]
この\bookterm{matrix}{行列}は,二つの値$x,y$を受け取り,新しい二つの値$2x+y,x+2y$を作る規則を表している.

一般には
\[
  \begin{pmatrix}p&q\\r&s\end{pmatrix}
  \begin{pmatrix}x\\y\end{pmatrix}
  =\begin{pmatrix}px+qy\\rx+sy\end{pmatrix}
\]
と定める.
したがって,連立した二つの式
\[
  x'=px+qy,\qquad y'=rx+sy
\]
を,\bookterm{matrix}{行列}と\bookterm{column-vector}{列ベクトル}の積一つにまとめられる.
ここで$x',y'$は,更新後の値を表す記号である.

また,成分ごとに展開すれば
\[
  A(c\boldsymbol{u}+d\boldsymbol{v})
  =cA\boldsymbol{u}+dA\boldsymbol{v}
\]
が成り立つ.
これは,二つの\bookterm{column-vector}{列ベクトル}を足してから更新しても,別々に更新してから足しても結果が同じ,という性質である.
後で初期値を二つの部分に分けて考える際に使う.

\subsection{行列の積は,更新を続けて行うことを表す}

\bookterm{matrix}{行列}をもう一度掛けると,同じ規則で二回更新できる.
先ほどの例なら
\[
  \begin{pmatrix}1\\0\end{pmatrix}
  \ \longmapsto\
  \begin{pmatrix}2\\1\end{pmatrix}
  \ \longmapsto\
  \begin{pmatrix}5\\4\end{pmatrix}
\]
と進む.
この二回分の更新を一つにまとめた\bookterm{matrix}{行列}を$A^2$と書く.
実際に文字$x,y$で計算すると
\[
  \begin{aligned}
    A\left(A\begin{pmatrix}x\\y\end{pmatrix}\right)
    &=A\begin{pmatrix}2x+y\\x+2y\end{pmatrix}\\
    &=\begin{pmatrix}5x+4y\\4x+5y\end{pmatrix}.
  \end{aligned}
\]
したがって
\[
  A^2=\begin{pmatrix}5&4\\4&5\end{pmatrix}
\]
である.
各成分を単に二乗するのではなく,更新を二回行った結果に合わせて積を定めている.

異なる\bookterm{matrix}{行列}$A,B$でも,まず$B$,次に$A$を掛ける操作を$AB$と書く.
$B$の各列に$A$を掛けた結果を並べると,積$AB$が得られる.
具体的には
\[
  \begin{aligned}
    &\begin{pmatrix}p&q\\r&s\end{pmatrix}
     \begin{pmatrix}u&v\\w&z\end{pmatrix}\\
    &\qquad=
     \begin{pmatrix}
       pu+qw&pv+qz\\
       ru+sw&rv+sz
     \end{pmatrix}.
  \end{aligned}
\]
右の\bookterm{matrix}{行列}から先に作用するため,一般には$AB$と$BA$は一致しない.

何も変えない更新は
\[
  I=\begin{pmatrix}1&0\\0&1\end{pmatrix},
  \qquad
  I\begin{pmatrix}x\\y\end{pmatrix}
  =\begin{pmatrix}x\\y\end{pmatrix}
\]
で表される.
$I$を\textbf{\bookterm{identity}{単位行列}}という.
数の累乗で$a^0=1$とするのに対応して,\bookterm{matrix}{行列}では$A^0=I$と定める.

\subsection{固有ベクトルと固有値：更新が定数倍になる場合}
  \BookFigEigen


毎回\bookterm{matrix}{行列}の積を計算するのは手間がかかる.
そこで,\bookterm{matrix}{行列}を掛けても定数倍になるだけの\bookterm{column-vector}{列ベクトル}を探してみよう.
先ほどの$A$について
\[
  \begin{aligned}
    A\begin{pmatrix}1\\1\end{pmatrix}
      &=\begin{pmatrix}3\\3\end{pmatrix}
       =3\begin{pmatrix}1\\1\end{pmatrix},\\
    A\begin{pmatrix}1\\-1\end{pmatrix}
      &=\begin{pmatrix}1\\-1\end{pmatrix}
  \end{aligned}
\]
である.
第一の\bookterm{column-vector}{列ベクトル}は一回ごとに$3$倍され,第二の\bookterm{column-vector}{列ベクトル}は変わらない.
したがって$k$回更新した結果は
\[
  A^k\begin{pmatrix}1\\1\end{pmatrix}
    =3^k\begin{pmatrix}1\\1\end{pmatrix},
  \qquad
  A^k\begin{pmatrix}1\\-1\end{pmatrix}
    =\begin{pmatrix}1\\-1\end{pmatrix}
\]
とすぐに分かる.

一般に,零でない\bookterm{column-vector}{列ベクトル}$\boldsymbol{u}$が
\[
  A\boldsymbol{u}=\lambda\boldsymbol{u}
\]
を満たすとき,\,$\boldsymbol{u}$を$A$の\textbf{\bookterm{eigenvector}{固有ベクトル}},\,$\lambda$を対応する\textbf{\bookterm{eigenvalue}{固有値}}という.
\bookterm{eigenvalue}{固有値}は,そのベクトルに対して一回の更新が何倍になるかを表す数である.
零ベクトル$\boldsymbol{0}=\begin{pmatrix}0\\0\end{pmatrix}$は,どんな$\lambda$でもこの式を満たしてしまうため,\bookterm{eigenvector}{固有ベクトル}から除く.
負の\bookterm{eigenvalue}{固有値}では矢印の向きが反転し,\bookterm{eigenvalue}{固有値}$0$では零ベクトルになる場合もある.

\subsubsection{ベクトルを二つの部分に分ける}

\bookterm{eigenvector}{固有ベクトル}が見つかると,ほかのベクトルへの作用も計算しやすくなる.
たとえば
\[
  \begin{pmatrix}1\\0\end{pmatrix}
  =\frac12\begin{pmatrix}1\\1\end{pmatrix}
   +\frac12\begin{pmatrix}1\\-1\end{pmatrix}
\]
である.
右辺を実際に足せば,上が$1$,下が$0$になる.
\sectionref*{節：一次独立・基底・次元}では,独立な解の\bookterm{combination}{線形結合}で\bookterm{initial}{初期条件}を合わせた.
ここでは二つの独立な\bookterm{column-vector}{列ベクトル}を\bookterm{basis}{基底}に選び,元のベクトルを同じように\bookterm{combination}{線形結合}で表している.
\bookterm{sequence}{数列}全体を分けた見方を,今度は更新前の状態の分解に使っているのである.
二つの部分を別々に$k$回更新すれば
\[
  \begin{aligned}
    A^k\begin{pmatrix}1\\0\end{pmatrix}
      &=\frac{3^k}{2}\begin{pmatrix}1\\1\end{pmatrix}
        +\frac12\begin{pmatrix}1\\-1\end{pmatrix}\\
      &=\begin{pmatrix}\frac{3^k+1}{2}\\\frac{3^k-1}{2}\end{pmatrix}.
  \end{aligned}
\]
大きな累乗を一回ずつ計算する代わりに,定数倍になる二つの部分を追えばよい.
本編で等比型に帰着させると計算が簡単になったのと同じように,\bookterm{matrix}{行列}でも一定の倍率で変わる部分を見つけることが鍵になる.

一般に,二つの異なる\bookterm{eigenvalue}{固有値}$\lambda_1,\lambda_2$に対応する\bookterm{eigenvector}{固有ベクトル}$\boldsymbol{u}_1,\boldsymbol{u}_2$があれば,どんな二成分のベクトルも
\[
  \boldsymbol{v}=c\boldsymbol{u}_1+d\boldsymbol{u}_2
\]
と表せる.
異なる\bookterm{eigenvalue}{固有値}に対応する二つの\bookterm{eigenvector}{固有ベクトル}は,互いの定数倍にならないからである.
係数$c,d$は,上下の成分を比較した連立方程式で求められる.
その後は
\[
  A^k\boldsymbol{v}
  =c\lambda_1^k\boldsymbol{u}_1+d\lambda_2^k\boldsymbol{u}_2
\]
で計算できる.
$k=0$のときは元の分解式を読み,\,$k\ge1$で\bookterm{eigenvalue}{固有値}が$0$なら対応する部分は零になる.

\subsection{行ベクトルと左固有ベクトル}

数を横に並べた$L=(\alpha,\beta)$を\textbf{\bookterm{row-vector}{行ベクトル}}という.
\bookterm{column-vector}{列ベクトル}$\boldsymbol{\ell}$を横に並べ直す操作を\textbf{\bookterm{transpose}{転置}}といい,\,$\boldsymbol{\ell}^{\mathsf T}$と書く.
たとえば$\boldsymbol{\ell}=\begin{pmatrix}\alpha\\\beta\end{pmatrix}$なら,\,$\boldsymbol{\ell}^{\mathsf T}=(\alpha,\beta)$である.
\bookterm{row-vector}{行ベクトル}と\bookterm{column-vector}{列ベクトル}の積は,対応する成分を掛けて足す.
\[
  (\alpha,\beta)\begin{pmatrix}x\\y\end{pmatrix}
  =\alpha x+\beta y.
\]
\bookterm{row-vector}{行ベクトル}を\bookterm{matrix}{行列}に掛けるときは,\bookterm{matrix}{行列}の各列との積を横に並べる.
たとえば,先ほどの$A$について
\[
  \begin{aligned}
    (1,\,1)A&=(3,\,3)=3(1,\,1),\\
    (1,-1)A&=(1,-1)
  \end{aligned}
\]
である.
一般に,零でない\bookterm{row-vector}{行ベクトル}$L$が$LA=\gamma L$を満たすとき,\,$L$を\textbf{\bookterm{left-eigen}{左固有ベクトル}}という.
この場合,\bookterm{column-vector}{列ベクトル}$\boldsymbol{v}$から作った数$L\boldsymbol{v}$は,更新後には
\[
  L(A\boldsymbol{v})=(LA)\boldsymbol{v}
  =\gamma L\boldsymbol{v}
\]
となる.
たとえば$L=(1,\,1)$なら,\,$x+y$という和が更新によって$3$倍される.
\bookterm{eigenvector}{固有ベクトル}は定数倍になる\bookterm{column-vector}{列ベクトル}を表し,\bookterm{left-eigen}{左固有ベクトル}は定数倍になる量を取り出す係数を表す.

\subsection{固有値を求める条件：行列式}

\bookterm{eigenvector}{固有ベクトル}の成分を比較すれば,\bookterm{eigenvalue}{固有値}を求める条件が得られる.
その条件を一つの式にまとめるため,\bookterm{matrix}{行列}式という記号を導入しよう.
二行二列の場合には
\[
  \det\begin{pmatrix}a&b\\c&d\end{pmatrix}=ad-bc
\]
と定める.
\bookterm{matrix}{行列}そのものではなく,四つの成分から計算した一つの数である.

二つの式$ax+by=0,\ cx+dy=0$について,\,$ad-bc\ne0$なら解は$x=y=0$だけになる.
たとえば,第一の式を$d$倍したものから第二の式を$b$倍したものを引くと,\,$(ad-bc)x=0$となる.
同様に$(ad-bc)y=0$も得られる.
逆に$ad-bc=0$なら,零でない解が存在する.
実際,\,$(a,b)\ne(0,\,0)$なら$(x,y)=(b,-a)$が両方の式を満たす.
$(a,b)=(0,\,0)$の場合も,残る一つの式から零でない解を選べる.

\bookterm{eigenvalue}{固有値}の式$A\boldsymbol{u}=\lambda\boldsymbol{u}$は
\[
  (\lambda I-A)\boldsymbol{u}=\boldsymbol{0}
\]
と書き直せる.
\bookterm{matrix}{行列}の引き算は対応する成分ごとに行うので,\,$A=\begin{pmatrix}p&q\\r&s\end{pmatrix}$なら
\[
  \lambda I-A
  =\begin{pmatrix}\lambda-p&-q\\-r&\lambda-s\end{pmatrix}.
\]
零でない解が存在する条件は
\[
  \begin{aligned}
    \det(\lambda I-A)
      &=(\lambda-p)(\lambda-s)-qr\\
      &=\lambda^2-(p+s)\lambda+(ps-qr)=0.
  \end{aligned}
\]
これを\bookterm{matrix}{行列}$A$の\textbf{\bookterm{characteristic}{特性方程式}}という.
この方程式を解いて$\lambda$を求め,成分の連立方程式を解けば,対応する\bookterm{eigenvector}{固有ベクトル}も見つかる.
たとえば$A=\begin{pmatrix}2&1\\1&2\end{pmatrix}$では
\[
  \lambda^2-4\lambda+3
  =(\lambda-3)(\lambda-1)=0
\]
だから,\bookterm{eigenvalue}{固有値}は$3,\,1$である.
これは数値例で確かめた倍率と一致する.

\bookterm{recurrence}{漸化式}では,次の項を計算するために記録しておく値の組を\textbf{状態}と呼ぶ.
また,\bookterm{matrix}{行列}式が$0$でないときは,更新後の二つの値から更新前の二つの値を一意に求め直せる.
このように更新を逆向きに戻せる\bookterm{matrix}{行列}を\textbf{可逆な\bookterm{matrix}{行列}}という.

次節では,これらの計算を本編の\bookterm{recurrence}{漸化式}に当てはめていこう.
特に,本編で選んだ「何を引くか」「どのような比や和を作るか」が,\bookterm{matrix}{行列}では何を選ぶことに対応するのかを考える.

% 付録・第1部：行列の基本から漸化式へ

\section{行列の基本から,特性方程式を見直す}\label{節：特性方程式を行列で見直す}

本節では,まず行列の読み方と計算の仕方を説明する.
その後,本編で使った「等比数列になるように置き換える」という工夫を,行列の言葉で見直していこう.
行列を初めて学ぶ読者も,最初の数値例から順に確かめてほしい.
ここでは主に二行二列の行列を扱う.

\subsection{行列とは何か}

行列とは,数を長方形に並べ,全体を括弧で囲んだものである.
たとえば
\[
  A=\begin{pmatrix}2&1\\1&2\end{pmatrix}
\]
は,縦に二段,横に二列の数を並べた行列である.
横の並びを\textbf{行},縦の並びを\textbf{列}という.
上が第一行,下が第二行であり,左が第一列,右が第二列である.
並んでいる一つ一つの数を\textbf{成分}と呼ぶ.
この例の第一行・第二列の成分は$1$である.

また,二つの数を縦に並べた
\[
  \boldsymbol{v}=\begin{pmatrix}x\\y\end{pmatrix}
\]
を\textbf{列ベクトル}と呼ぶ.
太字の$\boldsymbol{v}$は,二つの数をまとめて表すための記号である.
平面上の点$(x,y)$への矢印と考えることもできるが,以下ではまず「二つの値を入れた一組の記録」として読めばよい.

列ベクトルの足し算と定数倍は,成分ごとに行う.
\[
  \begin{pmatrix}x\\y\end{pmatrix}
  +\begin{pmatrix}u\\v\end{pmatrix}
  =\begin{pmatrix}x+u\\y+v\end{pmatrix},
  \qquad
  c\begin{pmatrix}x\\y\end{pmatrix}
  =\begin{pmatrix}cx\\cy\end{pmatrix}.
\]
たとえば$\begin{pmatrix}3\\1\end{pmatrix}+\begin{pmatrix}1\\-1\end{pmatrix}
=\begin{pmatrix}4\\0\end{pmatrix}$である.
二つの列ベクトルが等しいとは,上下それぞれの成分が等しいことをいう.

\subsubsection{行列を列ベクトルに掛ける}

先ほどの行列$A$を列ベクトルに掛けるときは,次のように計算する.
\[
  \begin{pmatrix}2&1\\1&2\end{pmatrix}
  \begin{pmatrix}x\\y\end{pmatrix}
  =
  \begin{pmatrix}2x+y\\x+2y\end{pmatrix}.
\]
上の成分は,第一行の$2,\,1$を$x,y$にそれぞれ掛けて足したもの.
下の成分は,第二行の$1,\,2$を$x,y$にそれぞれ掛けて足したものである.
たとえば
\[
  A\begin{pmatrix}1\\0\end{pmatrix}
  =\begin{pmatrix}2\\1\end{pmatrix},
  \qquad
  A\begin{pmatrix}2\\1\end{pmatrix}
  =\begin{pmatrix}5\\4\end{pmatrix}.
\]
この行列は,二つの値$x,y$を受け取り,新しい二つの値$2x+y,x+2y$を作る規則を表している.

一般には
\[
  \begin{pmatrix}p&q\\r&s\end{pmatrix}
  \begin{pmatrix}x\\y\end{pmatrix}
  =\begin{pmatrix}px+qy\\rx+sy\end{pmatrix}
\]
と定める.
したがって,連立した二つの式
\[
  x'=px+qy,\qquad y'=rx+sy
\]
を,行列と列ベクトルの積一つにまとめられる.
ここで$x',y'$は,更新後の値を表す記号である.

また,成分ごとに展開すれば
\[
  A(c\boldsymbol{u}+d\boldsymbol{v})
  =cA\boldsymbol{u}+dA\boldsymbol{v}
\]
が成り立つ.
これは,二つの列ベクトルを足してから更新しても,別々に更新してから足しても結果が同じ,という性質である.
後で初期値を二つの部分に分けて考える際に使う.

\subsection{行列の積は,更新を続けて行うことを表す}

行列をもう一度掛けると,同じ規則で二回更新できる.
先ほどの例なら
\[
  \begin{pmatrix}1\\0\end{pmatrix}
  \ \longmapsto\
  \begin{pmatrix}2\\1\end{pmatrix}
  \ \longmapsto\
  \begin{pmatrix}5\\4\end{pmatrix}
\]
と進む.
この二回分の更新を一つにまとめた行列を$A^2$と書く.
実際に文字$x,y$で計算すると
\[
  \begin{aligned}
    A\left(A\begin{pmatrix}x\\y\end{pmatrix}\right)
    &=A\begin{pmatrix}2x+y\\x+2y\end{pmatrix}\\
    &=\begin{pmatrix}5x+4y\\4x+5y\end{pmatrix}.
  \end{aligned}
\]
したがって
\[
  A^2=\begin{pmatrix}5&4\\4&5\end{pmatrix}
\]
である.
各成分を単に二乗するのではなく,更新を二回行った結果に合わせて積を定めている.

異なる行列$A,B$でも,まず$B$,次に$A$を掛ける操作を$AB$と書く.
$B$の各列に$A$を掛けた結果を並べると,積$AB$が得られる.
具体的には
\[
  \begin{aligned}
    &\begin{pmatrix}p&q\\r&s\end{pmatrix}
     \begin{pmatrix}u&v\\w&z\end{pmatrix}\\
    &\qquad=
     \begin{pmatrix}
       pu+qw&pv+qz\\
       ru+sw&rv+sz
     \end{pmatrix}.
  \end{aligned}
\]
右の行列から先に作用するため,一般には$AB$と$BA$は一致しない.

何も変えない更新は
\[
  I=\begin{pmatrix}1&0\\0&1\end{pmatrix},
  \qquad
  I\begin{pmatrix}x\\y\end{pmatrix}
  =\begin{pmatrix}x\\y\end{pmatrix}
\]
で表される.
$I$を\textbf{単位行列}という.
数の累乗で$a^0=1$とするのに対応して,行列では$A^0=I$と定める.

ここで,二つの数列が
\[
  x_{n+1}=2x_n+y_n,\qquad
  y_{n+1}=x_n+2y_n
\]
を満たす場合を考えよう.
$n$番目の二つの値を
\[
  \boldsymbol{v}_n=\begin{pmatrix}x_n\\y_n\end{pmatrix}
\]
とまとめれば,\,$\boldsymbol{v}_{n+1}=A\boldsymbol{v}_n$である.
このように,次の値を計算するために持っておく値の組を\textbf{状態}と呼ぶ.
初期状態から一回,二回と進めると
\[
  \boldsymbol{v}_2=A\boldsymbol{v}_1,\qquad
  \boldsymbol{v}_3=A^2\boldsymbol{v}_1
\]
なので,一般に
\[
  \boldsymbol{v}_n=A^{n-1}\boldsymbol{v}_1
\]
となる.
初項から第$n$項までの更新は$n-1$回である.
行列の累乗が分かれば,連立漸化式の一般項も分かることになる.

\subsection{固有ベクトルと固有値：更新が定数倍になる場合}

毎回行列の積を計算するのは手間がかかる.
そこで,行列を掛けても定数倍になるだけの列ベクトルを探してみよう.
先ほどの$A$について
\[
  \begin{aligned}
    A\begin{pmatrix}1\\1\end{pmatrix}
      &=\begin{pmatrix}3\\3\end{pmatrix}
       =3\begin{pmatrix}1\\1\end{pmatrix},\\
    A\begin{pmatrix}1\\-1\end{pmatrix}
      &=\begin{pmatrix}1\\-1\end{pmatrix}
  \end{aligned}
\]
である.
第一の列ベクトルは一回ごとに$3$倍され,第二の列ベクトルは変わらない.
したがって$k$回更新した結果は
\[
  A^k\begin{pmatrix}1\\1\end{pmatrix}
    =3^k\begin{pmatrix}1\\1\end{pmatrix},
  \qquad
  A^k\begin{pmatrix}1\\-1\end{pmatrix}
    =\begin{pmatrix}1\\-1\end{pmatrix}
\]
とすぐに分かる.

一般に,零でない列ベクトル$\boldsymbol{u}$が
\[
  A\boldsymbol{u}=\lambda\boldsymbol{u}
\]
を満たすとき,\,$\boldsymbol{u}$を$A$の\textbf{固有ベクトル},\,$\lambda$を対応する\textbf{固有値}という.
固有値は,そのベクトルに対して一回の更新が何倍になるかを表す数である.
零ベクトル$\boldsymbol{0}=\begin{pmatrix}0\\0\end{pmatrix}$は,どんな$\lambda$でもこの式を満たしてしまうため,固有ベクトルから除く.
負の固有値では矢印の向きが反転し,固有値$0$では零ベクトルになる場合もある.

\subsubsection{初期状態を二つに分けると,一般項が出る}

初期値を$(x_1,y_1)=(1,\,0)$として,先ほどの連立漸化式を解いてみよう.
初期状態は
\[
  \begin{pmatrix}1\\0\end{pmatrix}
  =\frac12\begin{pmatrix}1\\1\end{pmatrix}
   +\frac12\begin{pmatrix}1\\-1\end{pmatrix}
\]
と分けられる.
右辺を実際に足せば,上が$1$,下が$0$になる.
二つの部分を別々に$n-1$回更新すると
\[
  \begin{pmatrix}x_n\\y_n\end{pmatrix}
  =\frac{3^{n-1}}2\begin{pmatrix}1\\1\end{pmatrix}
   +\frac12\begin{pmatrix}1\\-1\end{pmatrix}.
\]
上下の成分を読み取って
\[
  x_n=\frac{3^{n-1}+1}{2},\qquad
  y_n=\frac{3^{n-1}-1}{2}
\]
を得る.
初期状態を,更新の簡単な二つの部分に分けたことが解法の要点である.

本編の連立型では,和と差を
\[
  u_n=x_n+y_n,\qquad v_n=x_n-y_n
\]
と置いた.
もとの二つの式を足す,引くという計算から
\[
  u_{n+1}=3u_n,\qquad v_{n+1}=v_n
\]
となる.
そして
\[
  x_n=\frac{u_n+v_n}{2},\qquad
  y_n=\frac{u_n-v_n}{2}
\]
で元に戻せる.
ここに現れた倍率$3,\,1$は,先ほどの固有値と同じである.

和と差を作る係数を横に並べると,\,$(1,\,1)$と$(1,-1)$になる.
このように数を横に並べたものを\textbf{行ベクトル}という.
行ベクトルと列ベクトルの積は,対応する成分を掛けて足す.
\[
  (\alpha,\beta)\begin{pmatrix}x\\y\end{pmatrix}
  =\alpha x+\beta y.
\]
行ベクトルを行列に掛けるときは,行列の各列との積を横に並べる.
たとえば$(1,\,1)A=(3,\,3)=3(1,\,1)$である.
一般の行列$A$でも,零でない行ベクトル$L$が
\[
  LA=\gamma L
\]
を満たせば,\,$c_n=L\boldsymbol{v}_n$は
\[
  c_{n+1}=LA\boldsymbol{v}_n=\gamma c_n
\]
を満たす.
この$L$を\textbf{左固有ベクトル}と呼ぶ.
本編で等比数列になるように$\alpha x_n+\beta y_n$の係数を選ぶ方法は,この行ベクトルを探す計算として理解できる.
固有ベクトルで状態を分ける方法と,等比数列になる量を取り出す方法は,同じ更新を異なる側から見ている.

\subsection{本編の特性方程式型を,行列で書き直す}\label{節：一つの値の更新を行列で書く}

次に,本編で扱った
\[
  a_{n+1}=pa_n+q\qquad(p\ne1)
\]
に戻ろう.
ここでは一つの値だけを更新しているが,定数$q$の足し算がある.
この足し算を行列の掛け算の中に含めるため,\,$a_n$と一緒に$1$も並べておく.
すると
\[
  \begin{pmatrix}p&q\\0&1\end{pmatrix}
  \begin{pmatrix}a_n\\1\end{pmatrix}
  =\begin{pmatrix}pa_n+q\\1\end{pmatrix}
  =\begin{pmatrix}a_{n+1}\\1\end{pmatrix}.
\]
上の成分では$q$に$1$を掛けることで,必要な足し算を作っている.
下の成分は毎回$1$のままなので,次の更新でも同じ計算ができる.

この行列を$H$とし,本編と同じく
\[
  \alpha=p\alpha+q,\qquad
  \alpha=\frac{q}{1-p}
\]
を満たす$\alpha$を取る.
この値から始めると次の項も変わらないので,\,$\alpha$は不動点である.
行列で計算すると
\[
  H\begin{pmatrix}\alpha\\1\end{pmatrix}
  =\begin{pmatrix}\alpha\\1\end{pmatrix},
  \qquad
  H\begin{pmatrix}1\\0\end{pmatrix}
  =p\begin{pmatrix}1\\0\end{pmatrix}.
\]
これで,更新が簡単な二つの列ベクトルが見つかった.

初期状態を
\[
  \begin{pmatrix}a_1\\1\end{pmatrix}
  =\begin{pmatrix}\alpha\\1\end{pmatrix}
   +(a_1-\alpha)\begin{pmatrix}1\\0\end{pmatrix}
\]
と分ける.
第一の部分は変わらず,第二の部分だけが一回ごとに$p$倍されるから
\[
  \begin{pmatrix}a_n\\1\end{pmatrix}
  =\begin{pmatrix}\alpha\\1\end{pmatrix}
   +(a_1-\alpha)p^{n-1}\begin{pmatrix}1\\0\end{pmatrix}.
\]
上の成分を読めば
\[
  a_n=\alpha+(a_1-\alpha)p^{n-1}
\]
となる.
$p=0$では$a_2$以降が$\alpha=q$になるので直接読めばよい.

たとえば$a_{n+1}=2a_n+1,\ a_1=0$なら$\alpha=-1$である.
初期状態は
\[
  \begin{pmatrix}0\\1\end{pmatrix}
  =\begin{pmatrix}-1\\1\end{pmatrix}
   +\begin{pmatrix}1\\0\end{pmatrix}
\]
と分かれ,後半だけが$2$倍ずつになる.
したがって$a_n=-1+2^{n-1}$である.
本編の$b_n=a_n-\alpha$という置き換えは,この「不動点からのずれ」だけを取り出していたのである.

ここで用いた$\alpha=p\alpha+q$は,不動点を求める方程式である.
一方,行列の固有値を求める方程式も線形代数では「特性方程式」と呼ばれる.
二つは異なる式であり,不動点$\alpha$そのものを固有値と呼ぶわけではない.
上の行列$H$の固有値は$1$と$p$である.

$p=1$なら,もとの式は$a_{n+1}=a_n+q$である.
更新を$k$回続けると
\[
  H^k=\begin{pmatrix}1&kq\\0&1\end{pmatrix}
\]
となり,\,$a_n=a_1+(n-1)q$を得る.
$q\ne0$では不動点がなく,加えた$q$が更新回数だけ積み重なる.
これは後で見る,特性根が重なると一般項に$n$が現れる例につながっている.

\subsection{二つの項を記録すれば,三項間漸化式も行列になる}\label{節：状態を二つ持つと,二階の漸化式になる}

三項間漸化式では,次の項を計算するために直前の二つの項が必要である.
たとえばフィボナッチ数列
\[
  F_0=0,\quad F_1=1,\quad
  F_{n+2}=F_{n+1}+F_n\quad(n\ge0)
\]
を考える.
状態を$\begin{pmatrix}F_{n+1}\\F_n\end{pmatrix}$とすれば
\[
  \begin{pmatrix}1&1\\1&0\end{pmatrix}
  \begin{pmatrix}F_{n+1}\\F_n\end{pmatrix}
  =\begin{pmatrix}F_{n+1}+F_n\\F_{n+1}\end{pmatrix}
  =\begin{pmatrix}F_{n+2}\\F_{n+1}\end{pmatrix}.
\]
上の成分で新しい項を作り,下の成分に直前の項を残している.
実際,
\[
  \begin{pmatrix}1\\0\end{pmatrix}
  \longmapsto\begin{pmatrix}1\\1\end{pmatrix}
  \longmapsto\begin{pmatrix}2\\1\end{pmatrix}
  \longmapsto\begin{pmatrix}3\\2\end{pmatrix}
\]
と進む.

一般に
\[
  a_{n+2}=s a_{n+1}+t a_n\qquad(n\ge0)
\]
は
\[
  \boldsymbol{w}_n=\begin{pmatrix}a_{n+1}\\a_n\end{pmatrix},
  \qquad
  B=\begin{pmatrix}s&t\\1&0\end{pmatrix}
\]
と置けば,\,$\boldsymbol{w}_{n+1}=B\boldsymbol{w}_n$になる.
直前の二項を必要とするため,この形を\textbf{二階の漸化式}とも呼ぶ.
本編の$a_{n+2}+pa_{n+1}+qa_n=0$では,\,$s=-p,\ t=-q$とすればよい.

\subsubsection{なぜ本編と同じ特性方程式が出るのか}

$B$を掛けると$\lambda$倍になる,零でない列ベクトルを探そう.
成分を$x,y$と書けば
\[
  B\begin{pmatrix}x\\y\end{pmatrix}
  =\lambda\begin{pmatrix}x\\y\end{pmatrix}
  \quad\Longleftrightarrow\quad
  \begin{cases}
    sx+ty=\lambda x,\\
    x=\lambda y.
  \end{cases}
\]
二つ目の式から,\,$y=0$なら$x=0$になってしまう.
そこで$y\ne0$と分かるので,ベクトル全体を$y$で割り,\,$y=1,\ x=\lambda$としてよい.
一つ目の式は
\[
  s\lambda+t=\lambda^2,
  \qquad
  \lambda^2-s\lambda-t=0
\]
となる.
これが固有値を求める方程式である.

本編では,等比数列の形$a_n=\lambda^n$を漸化式に当てはめて,同じ方程式を得た.
行列では「二項をまとめた記録が一回ごとに$\lambda$倍になる条件」を調べている.
数列が等比的に進むことと,状態ベクトルが定数倍になることは,同じ条件を別の形で表している.
$\lambda=0$が根になる場合も,行列の方程式はそのまま使える.

二つの異なる根$\lambda_1,\lambda_2$があれば
\[
  \begin{pmatrix}\lambda_1\\1\end{pmatrix},
  \qquad
  \begin{pmatrix}\lambda_2\\1\end{pmatrix}
\]
が固有ベクトルになる.
初期状態を
\[
  \begin{pmatrix}a_1\\a_0\end{pmatrix}
  =c\begin{pmatrix}\lambda_1\\1\end{pmatrix}
   +d\begin{pmatrix}\lambda_2\\1\end{pmatrix}
\]
と分けるには,\,$c+d=a_0,\ c\lambda_1+d\lambda_2=a_1$を解けばよい.
根が異なるので,\,$c,d$は一意に決まる.
$n$回更新した後は,下の成分から
\[
  a_n=c\lambda_1^n+d\lambda_2^n
\]
を得る.
根が$0$の場合,\,$n=0$での係数は初期状態の式から読み,\,$n\ge1$ではその成分は$0$になる.

フィボナッチ数列なら特性方程式は$\lambda^2-\lambda-1=0$で,根は
\[
  \varphi=\frac{1+\sqrt5}{2},\qquad
  \psi=\frac{1-\sqrt5}{2}
\]
である.
初期状態$\begin{pmatrix}1\\0\end{pmatrix}$に合わせる条件は
\[
  c+d=0,\qquad c\varphi+d\psi=1.
\]
$\varphi-\psi=\sqrt5$より$c=1/\sqrt5,\ d=-1/\sqrt5$と決まり,
\[
  F_n=\frac{\varphi^n-\psi^n}{\sqrt5}
\]
を得る.
$|\psi|<1<\varphi$なので,\,$\psi^n$の成分は小さくなり,\,$\varphi^n$の成分が大きくなる.
固有値は一般項の計算に加えて,更新を重ねたときの増え方も教えている.
異なる二つの根が複素数の場合も,成分や係数を複素数まで広げれば同じ計算ができる.
実数の初期値に合わせると,共役な二つの根の寄与が組になり,数列の項は実数になる.

\subsubsection{行列式を使うと,固有値の条件をまとめられる}

ここまでのように成分を比較すれば,固有値は求められる.
その条件を一つの式にまとめる記号が\textbf{行列式}である.
二行二列の場合には
\[
  \det\begin{pmatrix}a&b\\c&d\end{pmatrix}=ad-bc
\]
と定める.
行列そのものではなく,四つの成分から計算した一つの数である.

二つの式$ax+by=0,\ cx+dy=0$について,\,$ad-bc\ne0$なら解は$x=y=0$だけになる.
たとえば,第一の式を$d$倍したものから第二の式を$b$倍したものを引くと,\,$(ad-bc)x=0$となる.
同様に$(ad-bc)y=0$も得られる.
逆に$ad-bc=0$なら,零でない解が存在する.
実際,\,$(a,b)\ne(0,\,0)$なら$(x,y)=(b,-a)$が両方の式を満たす.
$(a,b)=(0,\,0)$の場合も,残る一つの式から零でない解を選べる.

固有値の式$B\boldsymbol{u}=\lambda\boldsymbol{u}$は,
\[
  (\lambda I-B)\boldsymbol{u}=\boldsymbol{0}
\]
と書き直せる.
行列の引き算は対応する成分ごとに行うので
\[
  \lambda I-B
  =\begin{pmatrix}\lambda-s&-t\\-1&\lambda\end{pmatrix}.
\]
零でない解が存在する条件は
\[
  \begin{aligned}
    \det(\lambda I-B)
      &=(\lambda-s)\lambda-(-t)(-1)\\
      &=\lambda^2-s\lambda-t=0.
  \end{aligned}
\]
こうして,先ほど成分から求めた特性方程式が再び得られる.
一般の行列$A=\begin{pmatrix}p&q\\r&s\end{pmatrix}$でも,同じ計算で
\[
  \det(\lambda I-A)
  =\lambda^2-(p+s)\lambda+(ps-qr)=0
\]
が固有値を求める特性方程式になる.

\subsubsection{特性根が重なるとき}

特性方程式が重解$r\ne0$を持つときには,異なる二つの固有ベクトルに分ける先ほどの方法が使えない.
一般項に$n$が現れる理由は,もとの漸化式を直接変形すると分かりやすい.
重解のときは
\[
  a_{n+2}=2r a_{n+1}-r^2a_n
\]
なので,\,$a_n=r^n b_n$と置いて$r^{n+2}$で割ると
\[
  b_{n+2}-2b_{n+1}+b_n=0.
\]
これは
\[
  b_{n+2}-b_{n+1}=b_{n+1}-b_n
\]
ということだから,\,$b_n$は等差数列になる.
したがって
\[
  b_n=C+Dn,\qquad a_n=(C+Dn)r^n
\]
である.
等比的な倍率$r^n$を取り除くと,残りに等差的な増え方が現れるのである.
$r=0$が重解なら$a_{n+2}=0$なので,\,$a_2$以降はすべて$0$として直接扱う.

\subsection{一次分数型は,二つの成分の比として読む}

ここまで,列ベクトルの成分そのものを数列の項として読んだ.
今度は成分の\textbf{比}を読むと,本編の一次分数型につながることを見てみよう.
\[
  a_{n+1}=\frac{pa_n+q}{ra_n+s}
\]
に対し,\,$r\ne0,\ ps-qr\ne0$とし,すべての項で$ra_n+s\ne0$と仮定する.
初めを$(x_1,y_1)=(a_1,\,1)$として
\[
  \begin{pmatrix}x_{n+1}\\y_{n+1}\end{pmatrix}
  =M\begin{pmatrix}x_n\\y_n\end{pmatrix},
  \qquad M=\begin{pmatrix}p&q\\r&s\end{pmatrix}
\]
と更新する.
$a_n=x_n/y_n$であれば
\[
  \begin{aligned}
    \frac{x_{n+1}}{y_{n+1}}
    &=\frac{px_n+qy_n}{rx_n+sy_n}\\
    &=\frac{p(x_n/y_n)+q}{r(x_n/y_n)+s}
     =\frac{pa_n+q}{ra_n+s}.
  \end{aligned}
\]
また,\,$y_{n+1}=(ra_n+s)y_n$なので,分母が零にならないという仮定のもとで$y_n\ne0$が保たれる.
したがって,各回の成分比はもとの漸化式の項になる.
ベクトル全体を零でない数で何倍しても比は同じである.
このように成分比に注目して行列の作用を見る方法を,射影変換の見方という.

不動点$\alpha$については
\[
  p\alpha+q=\alpha(r\alpha+s)
\]
なので
\[
  M\begin{pmatrix}\alpha\\1\end{pmatrix}
  =(r\alpha+s)\begin{pmatrix}\alpha\\1\end{pmatrix}.
\]
つまり,不動点に対応する列ベクトルは固有ベクトルであり,固有値は$r\alpha+s$である.
$\alpha$自身が固有値なのではなく,ベクトルの二成分の比になっている.

異なる二つの不動点$\alpha,\beta$がある場合を考え,
\[
  \lambda_\alpha=r\alpha+s,\qquad
  \lambda_\beta=r\beta+s
\]
と置く.
各状態を
\[
  \begin{pmatrix}x_n\\y_n\end{pmatrix}
  =c_n\begin{pmatrix}\alpha\\1\end{pmatrix}
   +d_n\begin{pmatrix}\beta\\1\end{pmatrix}
\]
と分ければ
\[
  c_{n+1}=\lambda_\alpha c_n,\qquad
  d_{n+1}=\lambda_\beta d_n
\]
となる.
二つの係数が,別々の等比数列として進むのである.

本編の置き換えとの関係を,成分から確かめよう.
$x_n=\alpha c_n+\beta d_n,\ y_n=c_n+d_n$だから
\[
  x_n-\alpha y_n=(\beta-\alpha)d_n,
  \qquad
  x_n-\beta y_n=(\alpha-\beta)c_n.
\]
したがって
\[
  \frac{a_n-\alpha}{a_n-\beta}
  =\frac{x_n-\alpha y_n}{x_n-\beta y_n}
  =-\frac{d_n}{c_n}.
\]
以下では$a_1\ne\beta$とする.
このとき$c_1\ne0$であり,\,$ps-qr\ne0$より$\lambda_\alpha\ne0$なので,\,$c_n\ne0$が保たれる.
ここで$\lambda_\alpha=0$なら$M\begin{pmatrix}\alpha\\1\end{pmatrix}=\boldsymbol{0}$となり,行列式が零でないことに反する.
$a_1=\beta$の場合は,\,$a_n=\beta$という定数列として扱えばよい.

係数の更新規則より
\[
  \frac{a_{n+1}-\alpha}{a_{n+1}-\beta}
  =\frac{\lambda_\beta}{\lambda_\alpha}
   \frac{a_n-\alpha}{a_n-\beta}.
\]
本編で$b_n=(a_n-\alpha)/(a_n-\beta)$を等比数列にできたのは,この二つの係数の比を取り出していたからである.
不動点の方程式
\[
  r z^2+(s-p)z-q=0
\]
の解と係数の関係から$r(\alpha+\beta)=p-s$なので,
\[
  p-r\alpha=r\beta+s,\qquad
  p-r\beta=r\alpha+s
\]
となる.
よって公比は本編と同じ
\[
  \frac{\lambda_\beta}{\lambda_\alpha}
  =\frac{p-r\alpha}{p-r\beta}
\]
である.

たとえば$a_{n+1}=2a_n/(1-a_n)$では
\[
  M=\begin{pmatrix}2&0\\-1&1\end{pmatrix}
\]
で,不動点$0,-1$に対応する固有値はそれぞれ$1,\,2$になる.
$\alpha=0,\ \beta=-1$とすると
\[
  b_n=\frac{a_n}{a_n+1},\qquad b_{n+1}=2b_n
\]
である.
この$2$は,二つの固有値の比$2/1$から現れていた.

\subsection{三角関数型漸化式からチェビシェフ多項式へ}\label{節：三角関数型漸化式からチェビシェフ多項式へ}

特性根が複素数になると,数列には振動が現れる.
実係数の漸化式では複素根は共役な組で現れ,一般には$\rho>0,\ 0<\theta<\pi$の条件下で
\[
  \rho(\cos\theta+i\sin\theta),
  \qquad
  \rho(\cos\theta-i\sin\theta)
\]
と書ける.
特性方程式の係数を比べると,対応する漸化式は
\[
  a_{n+2}=2\rho\cos\theta\,a_{n+1}-\rho^2a_n
\]
である.
ド・モアブルの公式から,特性根の $n$ 乗は振幅 $\rho^n$ と角度 $n\theta$ を持つ.
そのため実数解は
\[
  a_n=\rho^n\{C\cos(n\theta)+D\sin(n\theta)\}
\]
という形になる.
$\rho>1$ なら振幅は大きくなり,$\,0<\rho<1$ なら小さくなる.

特に振幅が変わらない $\rho=1$ の場合,特性根は $\cos\theta\pm i\sin\theta$ であり,漸化式は
\[
  a_{n+2}=2\cos\theta\,a_{n+1}-a_n
\]
となる.
余弦と正弦の加法定理を使うと,$\,\cos(n\theta)$ も $\,\sin(n\theta)$ もこの同じ漸化式を満たす.
したがって,二つの初期値に合わせて定数 $C,D$ を選べば
\[
  a_n=C\cos(n\theta)+D\sin(n\theta)
\]
と表せる.
複素数の根は,ただ計算を複雑にするものではない.
実数の数列にひそむ振動と,その振幅の増減を記録している.

ここで $x=\cos\theta$ と書き,角度の整数倍を $x$ の多項式で表せないか考えてみよう.
余弦の加法定理
\[
  \cos((n+1)\theta)+\cos((n-1)\theta)
  =2\cos\theta\cos(n\theta)
\]
を変形すると
\[
  \cos((n+1)\theta)
  =2x\cos(n\theta)-\cos((n-1)\theta)
\]
となる.
右辺は,一つ前と二つ前の余弦から次の余弦を作っている.
そこで同じ規則を多項式の列に移し,
\begin{align*}
  &T_0(x)=1,\qquad T_1(x)=x,\\
  &T_{n+1}(x)=2xT_n(x)-T_{n-1}(x)
  \qquad(n\ge1)
\end{align*}

と定める.
最初の二つが多項式で,次も足し算と掛け算だけで作られるので,すべての $T_n(x)$ は多項式になる.
この列をチェビシェフ多項式という.

この定義から
\[
  T_2(x)=2x^2-1,\qquad
  T_3(x)=4x^3-3x
\]
が得られる.
これらは二倍角・三倍角の公式の右辺と同じ形をしている.
実際,
\[
  T_n(\cos\theta)=\cos(n\theta)
\]
が成り立つ.
導出を確かめよう.
$n=0,1$ では
\[
  T_0(\cos\theta)=1=\cos0,\qquad
  T_1(\cos\theta)=\cos\theta
\]
で一致する.
また余弦の加法定理から
\[
  \cos((n+1)\theta)
  =2\cos\theta\cos(n\theta)-\cos((n-1)\theta)
\]
である.
これは $T_{n+1}(\cos\theta)$ を作る漸化式と同じである.
したがって,二つの初めの値が一致し,次の値を作る規則も一致しているので,帰納法によりすべての $n$ で
\[
  T_n(\cos\theta)=\cos(n\theta)
\]
となる.
チェビシェフ多項式は,余弦の整数倍の角を,角度を使わずに $x=\cos\theta$ の多項式として記録している.

この関係は,本編で扱った三角関数型漸化式
\[
  a_{n+1}=2a_n^2-1
\]
にも戻ってくる.
初項が $-1\le a_1\le1$ なら,ある角 $\theta$ を選んで $a_1=\cos\theta$ と書ける.
二倍角の公式と $T_2(x)=2x^2-1$ から
\[
  a_{n+1}=T_2(a_n)
\]
であり,$\,a_n=\cos(2^{n-1}\theta)$ と予想できる.
この予想は初項で正しく,成り立つと仮定すれば
\[
  a_{n+1}
  =2\cos^2(2^{n-1}\theta)-1
  =\cos(2^n\theta)
\]
となるので,帰納法で確かめられる.
さらに $a_1=\cos\theta$ と $T_m(\cos\theta)=\cos(m\theta)$ を使えば
\[
  a_n=\cos(2^{n-1}\theta)
  =T_{2^{n-1}}(a_1)
\]
である.
たとえば本編の $a_1=\sqrt3/2=\cos(\pi/6)$ では
\[
  a_n=T_{2^{n-1}}\left(\frac{\sqrt3}{2}\right)
  =\cos\left(2^{n-1}\frac{\pi}{6}\right).
\]
つまり,本編の二次式による更新を何度も繰り返すことは,角度を二倍ずつ進めることと同じである.
その更新を多項式の言葉で書き直したものがチェビシェフ多項式だった.

より一般に,$\,T_m(T_k(x))=T_{mk}(x)$ という合成の関係も,同じ角度の式から分かる.
$x=\cos\theta$ とすれば
\[
  T_m(T_k(x))
  =T_m(\cos(k\theta))
  =\cos(mk\theta)
  =T_{mk}(x).
\]
この等式は $[-1,1]$ のすべての $x$ で成り立ち,両辺が多項式なので,多項式の恒等式として実数全体で成り立つ.
したがって $T_2$ を繰り返し合成する操作は,添字を二倍ずつ進める $T_2,T_4,T_8,\ldots$ にまとまる.

\subsection{非線形な更新も,不動点の近くでは線形に見える}\label{節：非線形な更新も,不動点の近くでは線形に見える}

ここまでは,漸化式を行列で表し,線形な更新の方向ごとの倍率を見てきた.
最後に,一次分数型のような非線形の更新規則でも,不動点の近くでは同じ見方ができることを確かめよう.
一項前から次が決まる漸化式
\[
  a_{n+1}=f(a_n)
\]
を考え,$\,\alpha$ を不動点 $f(\alpha)=\alpha$ とする.
そこからのずれを $e_n=a_n-\alpha$ と置くと
\[
  e_{n+1}=f(\alpha+e_n)-f(\alpha)
\]
である.
$f$ が $\alpha$ で微分可能なら,微分係数の定義から
\[
  \frac{f(\alpha+e)-f(\alpha)}{e}\longrightarrow f'(\alpha)
  \qquad(e\longrightarrow0)
\]
である.
つまり,ずれ $e$ が十分小さい範囲では
\[
  f(\alpha+e)-f(\alpha)\approx f'(\alpha)e
\]
となる.
曲線を不動点のすぐ近くまで拡大すると接線に見える,という一次近似を使ったのである.

この近似は,収束についても手がかりを与える.
もし $|f'(\alpha)|<1$ なら,$\,|f'(\alpha)|$ より大きく $1$ より小さい数 $c$ を選べる.
微分係数の定義により,$\,a_n$ が $\alpha$ に十分近ければ
\[
  |e_{n+1}|
  =|f(\alpha+e_n)-f(\alpha)|
  \le c|e_n|
\]
となる.
誤差は一歩ごとに少なくとも一定の割合で縮み,いったん十分近くに入ればその近くにとどまって $\alpha$ へ近づく.
したがって $|f'(\alpha)|<1$ は,不動点が近くの初期値を引き寄せることを示す.
先ほどの行列では固有値が各方向の倍率を表していた.
ここでは微分係数 $f'(\alpha)$ が,不動点近くでのずれの倍率を表している.
線形代数と微分は別々の話に見えるが,どちらも「一歩ごとの変化を倍率で測る」道具としてつながっている.

% 付録：階差型・総和型から微積分と近似計算へ

\section{一歩の変化を細かくすると,微積分になるか}\label{節：差分から始まる離散的な微積分}
  \BookFigDifference


前節では,本編で等比数列を作った置き換えを,行列の作用として捉え直した.
ここでは,\sectionref{節：階差型}の「差から元の数列を求める」という考え方を出発点にしよう.
その逆向きの操作は,\sectionref{節：総和型}で使った「和から差を取って一般項を取り出す」という操作である.
この二つを一緒に見ると,差分と和が微分と積分にどう対応するかが分かる.

以下では,最初の項を$a_0$と書くことがある.
数列$a_0,a_1,a_2,\ldots$の隣り合う項の差
\[
  \Delta a_n=a_{n+1}-a_n
\]
を一階差分という.
本編で使った$a_{n+1}-a_n$に,\,$\Delta$という記号と「差分」という名前を付けたのである.
たとえば$a_n=n^2$なら
\[
  \Delta a_n=2n+1,\qquad \Delta^2a_n=2.
\]
二次式から差を取ると一次式になり,もう一回差を取ると定数になる.
ここで$\Delta^2a_n$は,\,$\Delta a_{n+1}-\Delta a_n$の意味である.

本編の特性方程式型の別解でも,添字を一つずらして引くことで定数項$q$を消し,
\[
  \Delta a_{n+1}=p\,\Delta a_n
\]
という等比型を作った.
その計算も,差分によって扱いやすい量を取り出した例だった.

逆に,差分を足し合わせれば元の数列を求め直せる.
\[
  \begin{aligned}
    a_n-a_0
      &=(a_1-a_0)+\cdots+(a_n-a_{n-1})\\
      &=\sum_{k=0}^{n-1}\Delta a_k.
  \end{aligned}
\]
途中の項が打ち消され,最初と最後だけが残る.
本編の階差型で和を取った操作は,離散的な積分と見直せる.
総和型で$S_{n+1}-S_n=a_{n+1}$とした操作は,その逆である.

この見方は,\sectionref{節：階差等比複合型}で引く多項式$g(n)$を探したことにもつながる.
特に$p=1$なら,\,$g(n+1)-g(n)=f(n)$を満たす式を探す問題になる.
たとえば$S_n=\sum_{k=1}^n k^2$については
\[
  S_n-S_{n-1}=n^2,\qquad S_0=0.
\]
二次式の差を逆向きに戻すので,三次式
$S_n=An^3+Bn^2+Cn+D$を候補にする.
その差は
\[
  S_n-S_{n-1}
  =3An^2+(-3A+2B)n+(A-B+C).
\]
係数比較から$\displaystyle A=\frac{1}{3}$, $\displaystyle B=\frac{1}{2}$, $\displaystyle C=\frac{1}{6}$となり,\,$S_0=0$から$D=0$である.
この候補が同じ初期値と差の関係を満たすことにより
\[
  S_n=\frac{n(n+1)(2n+1)}6
\]
が確かめられる.
和の公式を覚えて使うだけでなく,「差を取ると目的の式に戻るものは何か」と考えて作ることができた.

ここからは微分と一次近似を使い,差分と連続的な変化の関係を見よう.
時刻$t_n=nh$の値を$a_n$とすると,一単位時間あたりの変化は
\[
  \frac{a_{n+1}-a_n}{h}
\]
である.
ある微分可能な関数$y$の値$a_n=y(t_n)$を記録したものなら,刻み幅$h$を小さくすると,この差分商は微分係数$y'(t_n)$に近づく.

逆に,微分方程式$y'=F(t,y)$を数値で解くには,今の点での傾きから
\[
  a_{n+1}=a_n+hF(t_n,a_n)
\]
と次の値を予測できる.
これをオイラー法という.
正確な連続解の値と一致するとは限らないが,接線の方向へ一歩ずつ進む漸化式で,曲線を近似している.

本編の等比型と比較しやすい例として,\,$y'=-y$, $y(0)=1$を考えよう.
連続解は$y(t)=e^{-t}$で,時間が進むと$0$に近づく.
オイラー法では
\[
  a_{n+1}=(1-h)a_n,\qquad a_0=1
\]
だから,\,$a_n=(1-h)^n$である.
本編の等比型で調べたように,\,$0$への収束条件は$|1-h|<1$, すなわち$0<h<2$になる.
$0<h<1$では正のまま減少するが,\,$1<h<2$では符号を交互に変えながら近づく.
さらに$h>2$では絶対値が増大する.
本当の解が減衰していても,刻み幅を大きくすると,近似計算の列はまったく違う動きになりうる.
本編で等比数列の公比を調べた知識は,数値計算が安定するかを判断する道具にもなる.

また,正の時刻$t$を$N$等分して$\displaystyle h=\frac{t}{N}$とすれば,時刻$t$における近似値は
\[
  a_N=\left(1-\frac{t}{N}\right)^N
  \longrightarrow e^{-t}.
\]
刻み幅を細かくすることと,近似値が連続解へ近づくことが,この例では基本的な指数関数の極限で確認できる.
一般の微分方程式でも同様の収束を期待するが,その保証には関数や刻み幅の条件を調べる必要がある.

次に,\sectionref{節：ニュートン法型}で扱った近似計算へ戻ろう.
本編では,接線と横軸の交点を次の値に選ぶことで
\[
  a_{n+1}=\frac12\left(a_n+\frac2{a_n}\right)
\]
を作り,\,$\sqrt2$への誤差がほぼ二乗されることを見た.
ここではその具体例から,なぜ一般のニュートン法でも誤差が二乗程度になるのかへ進む.

方程式$g(x)=0$に対するニュートン法の更新は
\[
  x_{n+1}=x_n-\frac{g(x_n)}{g'(x_n)}
\]
である.
$g(\alpha)=0$, $g'(\alpha)\ne0$とし,根$\alpha$の近くで$g$は二回連続微分可能とする.
本編のずれの置き換えと同じく,\,$e_n=x_n-\alpha$とおくと
\[
  e_{n+1}
  =\frac{g'(x_n)e_n-g(x_n)}{g'(x_n)}.
\]
分子は,接線による変化と実際の変化との差である.
二階微分$g''$の絶対値がこの近くで$M$以下なら
\[
  \begin{aligned}
    g'(x_n)e_n-g(x_n)
      &=\int_\alpha^{x_n}\{g'(x_n)-g'(t)\}\,dt,\\
    |g'(x_n)e_n-g(x_n)|
      &\le \frac{M}{2}|e_n|^2.
  \end{aligned}
\]
二つ目の式は,平均値の定理から
$|g'(x_n)-g'(t)|\le M|x_n-t|$として積分すれば得られる.
積分の向きが逆になる$x_n<\alpha$でも,絶対値を取れば同じ評価になる.
また,\,$g'(\alpha)\ne0$と連続性により,十分近くでは$|g'(x_n)|\ge m$となる正の定数$m$を選べる.
したがって
\[
  |e_{n+1}|\le\frac{M}{2m}|e_n|^2.
\]
誤差が十分小さければ,次の誤差はその二乗の定数倍以下になる.
いったん十分近くから始めれば近くにとどまり,急速に根へ近づくこともこの評価から分かる.
本編の$\sqrt2$の例で得た正確な式
\[
  e_{n+1}=\frac{e_n^2}{2a_n}
\]
は,この一般的な性質が特に見えやすい形で現れたものだった.

前節の不動点の見方でも確かめられる.
更新関数$\displaystyle N(x)=x-\frac{g(x)}{g'(x)}$について
\[
  N'(\alpha)
  =\left.\frac{g(x)g''(x)}{g'(x)^2}\right|_{x=\alpha}=0.
\]
つまり,根のところでは一次のずれの倍率が零になり,二次のずれが残る.
本編で等比型へ帰着させたときは一定倍率でずれを追ったが,ニュートン法ではずれが小さいほど,次に残る割合も小さくなるのである.

接線の傾きを毎回微分で求めるのが難しいなら,過去二点を結ぶ割線の傾きで代用できる.
その更新は
\[
  x_{n+1}
  =x_n-g(x_n)\frac{x_n-x_{n-1}}{g(x_n)-g(x_{n-1})}
\]
となり,割線法と呼ばれる.
分母が零でない範囲で用いる.
直前の二項を記録して次を作るという点は,前節の二成分の状態と共通している.
本編で与えられた漸化式を解いた経験は,このように近似計算の規則を作り,その収束を調べることにもつながる.

% 付録再編草案・第3部：数列全体と母関数

\section{数列全体を一つの対象にまとめると,何が見えるか}

ここまでは,次の項をどう作るか,隣り合う項の差から何が読めるかを見てきた.
しかし,漸化式が決めているのは一つの項だけではない.
初項から先の全項をつなぐ,数列全体の骨組みである.
その骨組みを一度に持ち運べる形へまとめたら,何が見えるだろう.

数列
\[
  a_0,a_1,a_2,a_3,\ldots
\]
に対して,項の番号を印にして
\[
  A(x)=a_0+a_1x+a_2x^2+a_3x^3+\cdots
  =\sum_{n=0}^{\infty}a_nx^n
\]
と並べる.
これを母関数という.
$x$ を一つ掛けるごとに項の番号が一つ増えるので,\,$x$ は「何番目の項か」を記録する目盛りのような役割を持つ.
最初は,これを収束する関数としてではなく,数列を係数として並べた形式的な記号として扱おう.
母関数の力は,無限個の項を一つに圧縮することよりも,数列の操作を代数の操作に言い換えられることにある.

たとえば $A(x)$ に $x$ を掛けると,
\[
  xA(x)=a_0x+a_1x^2+a_2x^3+\cdots
\]
となる.
数列で言えば,初めに $0$ を置いて,その後の項を一つ右へずらした形である.
また,二つの数列の項を
\[
  c_n=\sum_{k=0}^{n}a_kb_{n-k}
\]
のように組み合わせる操作は畳み込みと呼ばれ,母関数を掛けたときの $x^n$ の係数になる.
一つの対象を作る方法を二つに分ける,あるいは二つの部品を組み合わせるとき,この畳み込みが自然に現れる.
つまり母関数では,「一つ先へずらす」は $x$ 倍に,「二つの選択を組み合わせる」は積に翻訳される.

この翻訳をフィボナッチ数列で追ってみよう.
\[
  F_0=0,\quad F_1=1,\qquad F_{n+2}=F_{n+1}+F_n
\]
とし,\,$F(x)=\sum_{n=0}^{\infty}F_nx^n$ と置く.
左辺の初めの項を取り出すと,
\[
  F(x)-x=\sum_{n=0}^{\infty}F_{n+2}x^{n+2}
\]
である.
漸化式 $F_{n+2}=F_{n+1}+F_n$ を代入すると,右辺の二つの和はそれぞれ $xF(x)$ と $x^2F(x)$ になる.
したがって
\[
  F(x)-x=xF(x)+x^2F(x)
\]
であり,
\[
  F(x)=\frac{x}{1-x-x^2}
\]
を得る.
項を一つずつ計算する代わりに,数列全体が満たす一つの代数方程式を解いた.
漸化式は母関数の分母に姿を変えたのである.

ここで分母を特性方程式と比べてみる.
\[
  1-x-x^2=0
\]
に対して $r=1/x$ と置くと,\,$r^2-r-1=0$ となる.
つまり母関数の分母の零点は,特性方程式の根の逆数である.
これは偶然ではない.
特性方程式は「一歩ごとの倍率」を調べ,母関数は「全項を並べたときの係数関係」を調べている.
行列の固有値,特性根,母関数の分母は,同じ更新規則を違う入口から見たものだった.

フィボナッチ数列をもう少し遠くから眺めると,長期的な増え方も見えてくる.
特性根は
\[
  \varphi=\frac{1+\sqrt5}{2},
  \qquad
  \psi=\frac{1-\sqrt5}{2}
\]
であり,\,$|\psi|<1<\varphi$ である.
一般項
\[
  F_n=\frac{\varphi^n-\psi^n}{\sqrt5}
\]
では,\,$\psi^n$ の影響は次第に小さくなる.
したがって
\[
  F_n\sim\frac{\varphi^n}{\sqrt5}
\]
となる.
ここで $\sim$ は,二つの量の比が $1$ に近づくという意味である.
大きな $n$ でどの根が勝つかを見れば,正確な一般項を毎回計算しなくても成長の景色が分かる.
母関数の分母からも同じ特性根が取り出されるため,固有値・一般項・母関数・漸近的な増大は一つの話につながる.

母関数が本領を発揮するのは,項どうしの組合せが漸化式に現れるときである.
カタラン数は,たとえば括弧の付け方や二分木の形を数える数列で,
\[
  C_0=1,\qquad C_{n+1}=\sum_{k=0}^{n}C_kC_{n-k}
\]
を満たす.
全体を左側の大きさ $k$ と右側の大きさ $n-k$ に分けると,左の形と右の形を一つずつ選ぶので積 $C_kC_{n-k}$ が現れる.
すべての分け方を足すため,右辺は畳み込みになっている.
母関数 $C(x)=\sum_{n=0}^{\infty}C_nx^n$ を作ると,畳み込みは積に変わり,
\[
  C(x)=1+xC(x)^2
\]
となる.
項ごとには複雑に見えた再帰が,母関数についての二次方程式へ変わった.
これを解くと
\[
  C(x)=\frac{1-\sqrt{1-4x}}{2x}
\]
である.
もう一方の符号では $x=0$ の近くで値が発散してしまい,定数項 $C_0=1$ を持つ母関数にならないため,こちらの枝を選ぶ.
母関数を使うと,非線形な漸化式を力ずくで代数方程式へ移し,そこから項の情報を引き出せる.

この発想は,チェビシェフ多項式にも戻ってくる.
その漸化式に母関数を作ると,
\[
  \sum_{n=0}^{\infty}T_n(x)t^n
  =\frac{1-xt}{1-2xt+t^2}
\]
となる.
これは,漸化式に対応する項を一つずつ並べ,同じ次数の係数を整理すれば得られる.
三角関数で定義した多項式列を,今度は一つの有理関数にまとめている.
漸化式,三角関数,母関数は別々の定義ではなく,同じ列を引き出すための三つの窓になっている.
このように,一つの対象に複数の表現を持たせると,見えにくかった性質が別の窓から見えてくる.

さらに遠くには,「数列を満たす関数を作る」話も広がっている.
階乗は整数 $n$ で $(n+1)!=(n+1)n!$ を満たすが,階乗を整数以外でも使いたい場面がある.
その拡張の一つがガンマ関数で,正の $x$ に対して
\[
  \Gamma(x)=\int_0^\infty t^{x-1}e^{-t}\,dt
\]
と定められる.
部分積分をすると
\[
  \Gamma(x+1)=x\Gamma(x)
\]
が現れ,整数では $\Gamma(n+1)=n!$ となる.
数列の外へ関数を作ると,階乗が積分や確率に現れる新しい景色へつながる.
本節ではその理論を追わないが,チェビシェフ多項式と同じく,漸化式が関数を生み出す一つの例として見ておこう.

確率に目を向けると, 状態の移り変わりも漸化式で記録できる.
晴れ・雨の順に確率を並べた列ベクトルを $\boldsymbol p_n$ とし, 晴れの翌日に晴れる確率を $0.8$, 雨の翌日に晴れる確率を $0.3$ とすると,
\[
  \boldsymbol p_{n+1}
  =
  \begin{pmatrix}0.8&0.3\\0.2&0.7\end{pmatrix}
  \boldsymbol p_n
  =P\boldsymbol p_n
\]
と書ける.
各列の和は $1$ なので, どの状態から出発しても次の確率の合計は $1$ に保たれる.
長く更新したときの定常分布は $P\boldsymbol p=\boldsymbol p$ を満たし, どの初期状態からそこへ近づくかは固有値で調べられる.
線形代数で見た状態の更新が, 確率の予測にもそのまま現れている.

数論では, ユークリッドの互除法が余りを次々に作る更新として表せる.
$a_{n-1}$ を $a_n$ で割った商を $q_n$, 余りを $a_{n+1}$ とすれば,
\[
  a_{n-1}=q_na_n+a_{n+1},
  \qquad 0\le a_{n+1}<a_n
\]
となる.
正の余りは毎回小さくなるため, この計算は有限回で止まり, 最大公約数が求まる.
ここでは一般項を求めず, 更新のたびに小さくなるという性質からアルゴリズムの停止を示した.

同じ数論の見方として, 整数係数が一定の二階漸化式を $m$ で割った余りで考えてみよう.
状態 $(a_n,a_{n+1})$ の各成分は $0,1,\ldots,m-1$ のいずれかなので, 状態は高々 $m^2$ 通りしかない.
次の状態が今の状態だけで決まるなら, いつか同じ状態が現れ, その後は同じ並びを繰り返す.
したがって, 余りの列は周期的になる.
たとえばフィボナッチ数列を $2$ で割った余りは
\[
  0,1,1,0,1,1,0,1,1,\ldots
\]
と周期 $3$ で繰り返す.
実数としての増大や収束を調べるときと, 余りとしての周期を調べるときでは, 同じ更新規則から見えてくる性質が変わる.

一般項を簡単な式で書きにくい数列にも, 調べられることは残っている.
たとえば
\[
  a_0=0,\qquad a_{n+1}=a_n^2+1
\]
では $0,1,2,5,26,\ldots$ と増える.
$n\ge2$ では $a_n\ge2$ であり, $a_{n+1}\ge a_n^2$ だから, 帰納的に
\[
  a_n\ge 2^{\,2^{n-2}}
\]
となる.
一般項を求めなくても, 増加が非常に速いことは分かる.
また, 数列が単調で有界と分かれば収束し, 更新が $a_{n+1}=f(a_n)$ で $f$ が連続なら, 極限 $L$ は $L=f(L)$ を満たす.
単調性, 有界性, 不動点, 増加速度は, 一般項とは別の問いに答えてくれる.

本編では, 漸化式から一般項を取り出す方法を整理した.
付録では, 同じ更新規則を行列で表し, 一歩の差を細かく見て微積分へつなぎ, 全項を母関数にまとめた.
確率の定常分布, 整数の余りの周期, 一般項を使わない増加の評価も, 見たい性質に応じて更新規則を読み替えた例である.
漸化式を理解する道具は, 一般項を得るためだけでなく, その数列がどのように動き, 何を数え, どのような性質を保つかを知るために広がっている.


\section{演習の操作を別の表現で見直す}\label{節：演習の操作を別の表現で見直す}
  新しい道具を学んだ後に,すでに解いた問題へ戻ろう.
  ここでは\bookproblemref{practice-entrance}{3}{発展問3},
  \bookproblemref{practice-entrance}{4}{問4},
  \bookproblemref{practice-entrance}{9}{問9}で行った操作を,
  \bookterm{generating}{母関数}・回転・\bookterm{invariant}{不変量}という見方で確かめる.

  \subsection{畳み込み和が一定になる理由}
  発展問3の$c_0=1$, $c_{n+1}=\frac{2n+1}{2n+2}c_n$から,
  \[
    c_n=\frac1{4^n}\binom{2n}{n}
  \]
  を得た.その\bookterm{generating}{母関数}を$C(x)=\sum_{n\ge0}c_nx^n$とする.
  まず,級数の各項を微分する形式的な操作を使う.
  収束する関数の微分を仮定するのではなく,
  $x^n$を$nx^{n-1}$へ移して係数を並べ直す操作として読む.
  元の\bookterm{recurrence}{漸化式}を$(n+1)c_{n+1}=(n+\frac12)c_n$と書けば,
  \[
    (1-x)C'(x)=\frac12C(x)
  \]
  を得る.積の係数は有限和なので,係数を比較することで通常の積の微分法則も使える.
  $D(x)=C(x)^2=\sum_{n\ge0}d_nx^n$と置くと,
  \[
    (1-x)D'(x)=2C(x)(1-x)C'(x)=D(x).
  \]
  $x^n$の係数を比較すれば$(n+1)d_{n+1}-nd_n=d_n$,すなわち$d_{n+1}=d_n$である.
  $d_0=c_0^2=1$なので,
  \[
    C(x)^2=\frac1{1-x}=\sum_{n\ge0}x^n
  \]
  の$x^n$の係数から
  \[
    \sum_{k=0}^n c_kc_{n-k}=1
  \]
  を得る.問題の解法は有限和を変形し,この解法は\bookterm{sequence}{数列}全体の積へ移した.
  どちらも同じ\bookterm{convolution}{畳み込み}を扱っている.

  \subsection{中点の更新を,回転と面積から見直す}
  発展問4では,平行四辺形の各辺上で同じ比率$t$ $(0<t<1)$の分点を取り,
  新しい平行四辺形を作った.連続する辺ベクトルを$u,v$とする.
  頂点の更新を引き算すれば,
  \[
    u'=(1-t)u+tv,\qquad v'=-tu+(1-t)v
  \]
  である.従って辺の組を混ぜる\bookterm{matrix}{行列}は
  \[
    M=\begin{pmatrix}1-t&t\\-t&1-t\end{pmatrix}
  \]
  となる.\bookterm{matrix}{行列}式は$\rho=(1-t)^2+t^2$だから,
  平行四辺形の面積は一回ごとに$\rho$倍となる.
  \bookterm{eigenvalue}{固有値}は$1-t\pm it$であり,その絶対値の二乗も$\rho$である.

  正方形で,向きを反時計回りに取った辺を複素数として$v=iu$と書けるときは,
  \[
    u'=((1-t)+it)u
  \]
  である.$t=\frac12$なら倍率は$\frac1{\sqrt2}$,回転角は$\frac\pi4$になる.
  これが問題の中点更新の辺長と向きに一致する.
  \begin{bookpoint}
    \textbf{\bookterm{matrix}{行列}が混ぜる対象を確認する}\par
    一般の平行四辺形では$u,v$は直交する等長の辺ではない.
    辺の組を混ぜる\bookterm{matrix}{行列}を,そのまま平面上の相似変換と読むことはできない.
    面積の倍率の説明は一般の場合に使え,上の回転の説明には正方形という条件を使う.
  \end{bookpoint}

  \subsection{整数方程式を保つ更新を,不変量から見る}
  発展問9では,$u_0=3$, $u_1=3$, $u_{n+2}=3u_{n+1}-u_n$から,
  $(3,u_n,u_{n+1})$が$x^2+y^2+z^2=xyz$の解になることを示した.
  初期値では
  \[
    u_n^2+u_{n+1}^2+9=3u_nu_{n+1}
  \]
  が成り立つ.一般の$x,y$に対して
  \[
    y^2+(3y-x)^2-3y(3y-x)=x^2+y^2-3xy
  \]
  だから,左辺の差は更新の前後で変わらず$-9$に保たれる.
  その結果,$u_nu_{n+2}=u_{n+1}^2+9$となり,
  \[
    u_{n+2}=\frac{u_{n+1}^2+9}{u_n}
  \]
  とも書ける.正値性は$u_0=u_1=3$と,
  $u_2=6$以降の増加から確認できる.
  \sectionref*{節：不変量から線形の関係}の定数$1$を$9$へ変えた更新だった.
  元の問題では\bookterm{linear}{線形}の更新から保存を示し,本編の新しい例では保存から\bookterm{linear}{線形}の更新を取り出した.
  同じ関係を二つの向きから使っている.

  \begin{bookpoint}
    \textbf{別解から何を持ち帰るか}\par
    計算が短くなるだけでなく,足す意味・面積の意味・保存される関係が見える.
    新しい表現を使った後は,元の問題の条件と答えへ戻って対応を確かめよう.
  \end{bookpoint}


\end{bookcolumns}
